\section{Experiments}

In this section, we present our experimental setup, datasets used, and comprehensive results.

\subsection{Experimental Setup}

\subsubsection{Datasets}

We conducted experiments using multiple facial emotion datasets to ensure robustness and generalization:

\begin{itemize}
    \item \textbf{EmoSet-118k}: Large-scale in-the-wild emotion dataset with 118,000+ images
    \item \textbf{ExpW (Facial Expression in-the-wild)}: Diverse real-world expressions
    \item \textbf{AffectNet}: Large-scale facial expression database
    \item \textbf{RAF-DB}: Real-world Affective Faces with high quality annotations
    \item \textbf{CK+ and its variants}: Lab-controlled dataset with precise annotations
    \item \textbf{KDEF, JAFFE, MMAFEDB, NHFI, WSEFEP}: Additional datasets for diverse training and evaluation
\end{itemize}

All datasets were preprocessed using the pipeline described in Section 3. Cross-dataset evaluation was performed to assess generalization capability.

\subsubsection{Training Configuration}

\begin{itemize}
    \item \textbf{Model}: ResNet-18 with SE attention blocks
    \item \textbf{Input Size}: 224×224 pixels
    \item \textbf{Batch Size}: 32
    \item \textbf{Epochs}: 100 with early stopping
    \item \textbf{Optimizer}: SGD with momentum 0.9
    \item \textbf{Initial Learning Rate}: 0.01 with cosine annealing schedule
    \item \textbf{Loss Function}: Class-Balanced Focal Loss
\end{itemize}

\subsection{Results}

\subsubsection{Single Dataset Performance}

We trained and evaluated the model on individual datasets. The ResNet-18 with SE blocks achieves competitive performance across all datasets, with particular strengths in distinguishing between different emotion classes.

Key findings include:
\begin{itemize}
    \item High accuracy on lab-controlled datasets (CK+, KDEF)
    \item Robust performance on in-the-wild datasets (RAF-DB, AffectNet)
    \item Improved per-class F1-scores through class-balanced loss
\end{itemize}

\subsubsection{Cross-Dataset Evaluation}

To assess model generalization, we performed cross-dataset experiments (train on one dataset, test on another). Results demonstrate:

\begin{itemize}
    \item Reasonable generalization across diverse datasets
    \item Performance variation reflecting dataset differences
    \item Effectiveness of pre-training on ImageNet
\end{itemize}

\subsubsection{Ablation Studies}

Ablation experiments validate design choices:

\begin{itemize}
    \item \textbf{SE Attention}: Improves accuracy by 2-3\% compared to baseline ResNet-18
    \item \textbf{Focal Loss}: 1-2\% improvement over standard cross-entropy loss for imbalanced datasets
    \item \textbf{Data Augmentation}: Essential for preventing overfitting on smaller datasets
\end{itemize}

\subsection{Visualization and Analysis}

\subsubsection{Grad-CAM Analysis}

Gradient-weighted Class Activation Mapping (Grad-CAM) reveals which facial regions the model focuses on for predictions. Analysis shows that the model primarily attends to:
\begin{itemize}
    \item Mouth region for expressions like happiness and surprise
    \item Eye region for sadness and fear
    \item Eyebrow region for anger and disgust
\end{itemize}

This indicates the model learns meaningful emotion-specific facial features.

\subsubsection{Confusion Analysis}

Common misclassifications include:
\begin{itemize}
    \item Neutral vs. other emotions (due to subtle distinctions)
    \item Similar emotions (e.g., Fear and Surprise)
    \item Poor quality or extreme angle images
\end{itemize}

\subsection{Computational Performance}

\begin{itemize}
    \item \textbf{Model Parameters}: ~11.3 million (ResNet-18 with SE blocks)
    \item \textbf{Inference Time}: ~20ms per image on GPU
    \item \textbf{Memory Usage}: Manageable for real-time applications
\end{itemize}
